% main.tex — arXiv rendering of the physical-iso paper.
% Source of truth: ../paper.md. This file is the gr-qc submission skeleton.
%
% arXiv categories: gr-qc (primary); cross-lists: quant-ph, cs.LO
% Build: pdflatex main.tex && pdflatex main.tex   (or: tectonic main.tex)
% COMPILED in this container: pdflatex, zero errors, main.pdf (10 pages);
% re-run after any edit and check main.log.
\documentclass[11pt]{article}
\usepackage[margin=1.1in]{geometry}
\usepackage{amsmath}
\usepackage{guarded-syntax} % sorts, operators, claim labels, tikz styles (from the notation artifact)

\title{Conservation as Grammar:\\ No-Cloning and Unitarity as Structural Theorems\\ of a Guarded Linear Calculus,\\ with a Holographic Resolution Bound}
\author{derivation-of-care research program\thanks{Companion mechanization \& program artifacts cited as [CCS], [GS], [GL].}}
\date{}

\begin{document}
\maketitle

\begin{abstract}
In standard quantum mechanics, unitarity is a postulate and no-cloning is a
theorem about the formalism's linear maps. We exhibit a mechanized, sorted,
guarded linear calculus in which both statements are \emph{structural}
theorems: properties of which well-formed processes exist at all, prior to any
dynamics. No derivation duplicates a linear token
(\dcn{no_contraction_content}, machine-checked in Lean~4), and in the guarded
kernel no derivation consumes one into nothing
(\dcn{no_vanishing_restricted}). The two laws carry different strength ---
cloning is forbidden by the grammar alone, even where discard and duplication
remain available for other sorts, while conservation of content requires the
kernel's discard ban --- an asymmetry that mirrors the physical one between a
kinematic and a dynamical law. The unrestricted calculus is not a harmless
generalization: there, information loss is \emph{derived} (an explicit
construction is machine-checked) and the unguarded dynamics provably diverge.
We give a translation manual into the language of observables, tensor
networks, and structural holography; for the finitary fragment the
dictionary is backed by a semantic bridge --- a functor into the category
of completely positive maps, under which the kernel/outside boundary is the
no-discard/discard boundary and the loss process denotes the trace --- while
the remaining rows are readings: a dictionary, not an isomorphism, whose
functor mirrors the laws without deriving them.
Finally, if the guarded tick is \emph{spatialized} --- an assumption the
calculus does not make --- combining it with the Bekenstein--Hawking area law
yields a bulk voxel scale $L \approx (l_p^2 R)^{1/3} \approx 5\times
10^{-15}\,\mathrm{m}$ for the observable radius $R$: a projection, not a
theorem. The exponent belongs to the entropy bound; the guard fixes no scale.
\end{abstract}

\section{Introduction}

Unitarity is a postulate: time evolution is required to preserve transition
probabilities and to be reversible, and the symmetry structure of pure states
then forces the implementing maps to be unitary up to phase [Wig59].
No-cloning is a theorem [WZ82]: a copier of unknown states would have to act
nonlinearly on superpositions. This paper presents a formal system --- the
Care Calculus [CCS], a sorted linear logic with a guarded fixpoint, mechanized
in Lean~4 --- in which both facts are theorems of the same kind: structural
non-derivability results about which processes exist. Their difference in
strength is reproduced by the system's architecture: cloning is forbidden even
in the most permissive rule-set; conservation fails there, by an explicit
machine-checked construction, and is restored by exactly one restriction ---
the discard ban. Every mapping to physics in this paper is labeled
\interp{reading}, every machine-checked anchor \meas{voucher named}, and the
one experimental section \proj{named instruments}.

\section{The calculus and its guarded kernel}

Sorts: three linear resource types --- the information token $C$ (content),
spent exactly once in every derivation (\dcn{single_linear_use} \meas{}); the
inert write-only archive (trace), with no elimination rule; and the live
guarded stream, spent like content, not copied --- and a non-linear family:
duplicable sources (knower), derived knowledge, values. The
fixpoint is guarded by the later modality $\guard{}$ (one constructor per
interval; [Nak00, AM13, Bir12]) and fires only at all-non-linear contexts:
\dcn{no_fix_at_content} excludes a fixpoint at the content frame \meas{} ---
at any position in any content-bearing context
(\dcn{guard_inherited} \meas{}) --- and \dcn{selfapp_price} prices
self-reference exactly
\meas{}. The \emph{Guarded Kernel} (the \dcn{ArmedSetup}) removes two rules:
the discard (\texttt{NoWeak}) and the knower-synthesis ban
(\dcn{NoKnowerSynthesis}). Within it both conservation laws are absolute:
\dcn{no_vanishing_restricted} and \dcn{no_contraction_restricted}
\meas{}, on the counting models of [CCS] (in the contraction model, sound at
all seventeen rules with no side condition; in the kernel's vanishing proof,
sound at sixteen with the weak rule's obligation discharged by the discard
ban itself). The axiom footprint: the kernel's theorems sit inside
\texttt{propext} and \texttt{Quot.sound} alone, the audit allowlist is
Lean's standard three --- no admitted axioms, no \texttt{sorry} in the
restricted path. The kernel is
not vacuous: the archive is legal (\dcn{archival_witness} \meas{}), the
closed unit stands, the constants stand. What it forbids is exactly the
discard. And the loss process needs no detours: a \emph{cut-free}
$[C] \vdash \mathbf{1}$ exists (\dcn{SCClosedKiller.witnessTyped_cutFree}
\meas{}) --- the leak theorem turns any cut-free non-linear output at the
token frame into one (\dcn{SCClosedKiller.killerOf_cutFree} \meas{}) --- so
the atomic-head restriction excludes nothing (\dcn{SCClosedKiller.lamArm_false}
\meas{}) and the discard ban excludes
everything (the witness carries weak nodes;
\dcn{SCClosedKiller.witness_not_noWeak} \meas{}).

\subsection{Translation table (excerpt)}

Linear token, used once $\mapsto$ a non-duplicable state;
\dcn{no_contraction_content} $\mapsto$ no-cloning [WZ82];
\dcn{no_vanishing_restricted} $\mapsto$ information conservation; the
lete/pair exchange $\mapsto$ a non-unitary channel (information delivered to
a sink nothing reads); \dcn{ArmedSetup} $\mapsto$ the boundary conditions of
a stable sector; the duplicable context vs.\ the linear stream $\mapsto$
boundary vs.\ bulk, structurally. All mappings \interp{}.

\section{The two conservation correspondences}

\textbf{No-cloning holds everywhere.} No derivation at $[C]$ concludes
$C\ten C$ (\dcn{no_contraction_content} \meas{}); no closed cloner exists
(\dcn{no_closed_cloner} \meas{}); any conclusion of $C \ten C$ costs a
frame of budget ${\ge}2$ (\dcn{contraction_needs_two} \meas{}). These hold
\emph{unrestricted} --- discard present, duplication available at the
duplicable sorts. The impossibility is prior to maps entirely: a grammar
fact. \interp{} The reading that linear types enforce this is standard in
quantum programming languages [SV06, PRZ17], with the Hilbert-space theorem
mechanized in Isabelle [BHL20], the categorical form in [Abr09, AC04], and
a denotational conservation-of-unitarity in [ILV26]; what the calculus adds
is the theorem's \emph{kind} --- a non-derivability in a grammar that also
contains the discard. The schema form (no \emph{linear type} is contractible) is
FALSE --- measured: any closed inhabitant of $\tau$ is a second copy at
$[\tau]$ (\dcn{contractible_of_closed} \meas{}); the concrete contractible
types are the trace family ($\mathrm{trace}$, $\guard\,\mathrm{trace}$,
$\mathrm{trace}\ten\mathrm{trace}$) and the self-arrows
$\mathrm{arrow}\,C\,C$ and $\mathrm{arrow}\,C\,\mathbf{1}$ --- closed by the
adversarial extractor itself --- those five linearities kernel computations
(\dcn{witnesses_are_linear} \meas{}), the generic self-arrow family following
at every type (\dcn{arrow_self_contractible} \meas{}). The repair is
Wootters--Zurek's own clause as
grammar: no-cloning is the law of the \emph{linear} types with \emph{no
closed inhabitant} --- you cannot clone what you cannot prepare ---
stated in-tree as the classification \dcn{T5}, its preparability half a
theorem at every type (\dcn{closed_implies_clone} \meas{}), the linearity
gate necessary (the knower: no closed inhabitant yet contractible,
\dcn{knower_contractible} \meas{}); the settled cases span both sides of the
boundary: unclonable --- the carriers $C$ and stream
(\dcn{no_contraction_stream} \meas{}) --- clonable --- the trace family and
every self-arrow. The frontier beyond them: a single wall [CCS] ---
$\mathrm{arrow}\,C\,\mathrm{stream}$, reducing to $[C] \nvdash
\mathrm{stream}$, the companion's own community ask, its base case now
in-tree (\dcn{no_closed_stream} \meas{}) and its tower case discharged;
one statement closes the ask and the classification at once.

And the asymmetry is \emph{computed} between adjacent types: the vanishing
conduit's type $\mathrm{arrow}\,K\,\mathrm{knowledge}$ is non-linear ---
weakable, the discard reaches it --- while the stream conduit's type
$\mathrm{arrow}\,K\,\mathrm{stream}$ is linear, because every stream-needing
type is (\dcn{needsS\_linear} \meas{}), so the same weakening is structurally
blocked [CCS]: the two laws differ by a single \texttt{Bool}. The rival
derivation --- both laws from a single conservation assumption [HHSS05] ---
runs the deduction opposite to ours. And the
construction's autopsy shows the same gap: the vanishing channel runs to the
discard --- a step the kernel must ban --- while the stream-cloning
construction dies at its \emph{trigger} (conduit inhabited, nothing fires),
needing no kernel at all [CCS].

\textbf{Conservation needs the kernel.} In the unrestricted calculus a
derivation consumes $C$ into the trivial type:
\dcn{no_vanishing_unrestricted_refuted} \meas{}. Under the discard ban the
law is absolute (\dcn{no_vanishing_restricted},
\dcn{no_vanishing_armed} \meas{}); the token's single use is pinned
(\dcn{no_weakening_content} \meas{}).

\textbf{The asymmetry} mirrors the physical one: no-cloning follows from
linearity alone (kinematic), information conservation requires unitary
dynamics (dynamical) [WZ82]. The calculus reproduces both laws \emph{and the
gap between them}. The proof method --- a monotone budget every rule
respects --- is conserved-charge reasoning, with a second engine in-tree: the
flow count, a net current from frame to conclusion with generator fires as
its sources (its content instance: no clone without a mint,
\dcn{no_contraction_nogen} \meas{}) --- conserved charge and conserved
current, different theorems about the same law. \interp{}

\textbf{The archive.} Content may be consumed into the inert record type
(``reception without reduction''): the record is counted, nothing reads it
back --- no closed \texttt{trace}→\texttt{knowledge} function exists
(\dcn{past_cannot_teach} \meas{}), nor \texttt{trace}→\texttt{value}
(\dcn{past_cannot_value} \meas{}); ``store now, understand later'' is
untypable. Under the dictionary this is the Pati--Braunstein no-deleting
territory [PB00], with a split the physics does not usually draw: the
kernel forbids the discard, and the grammar makes even the \emph{legal}
archive permanently unread --- two absolutes where physics states one.
\interp{} And on the information paradox the archive takes a stance:
what crosses is neither destroyed nor retrievable --- recorded, counted,
inert --- the shape of the retained-information resolutions [Page93,
Haw05]. \interp{}

\textbf{The semantic bridge.} For the finitary fragment the dictionary is
backed by a construction: a functor into $\mathbf{CPM}(\mathbf{FdHilb})$
[Sel07], the dagger-compact category of completely positive maps. The
discard denotes the trace (available at every object --- the unrestricted
calculus is exactly the channel world); copying exists only along a chosen
classical structure [AC04], a natural one collapsing the category [Abr09];
and the two laws are mirrored at the clause level: the only clause that
erases is the discard (any CP map $\mathcal{L}(H) \to \mathbb{C}$ is
$X \mapsto \mathrm{Tr}(AX)$ with $A \geq 0$, hence
$\mathrm{Tr} \circ \Phi_A$) --- the
semantic image of
\dcn{counterexample_not_noweak} --- while the cloner \dcn{no_contraction_content}
forbids would be a structureless copier, which cannot exist. The
interpretation covers the finitary fragment the table rows name: the sorts,
the four constants (fixed channels --- the coalgebra's faces), the archive
(\dcn{nil}/\dcn{snoc} along the classical structure), first-order arrows,
and the structural rules at their gated sorts. A grain note, found by
building the companion simulator [CCS-sim]: the calculus's own discard is
\emph{gated} --- weakening at non-linear types only --- so erasure of the
token requires first carrying it into a discardable sector: the leak
theorem's mechanism, denotationally. The
preparability law reads: you cannot copy along a structure you
have not chosen. Three limits: soundness does not reflect (the grammar's
no-go theorems are stronger than their denotations); the scope is finitary
and first-order; nothing here is mechanized. The guarded half --- a quantum
analogue of the topos of trees [Bir12] (the original is Set-valued) with
$\guard$ the shift, the fixpoint the productive
solution, one channel per tick --- is stated as a program, not a result.

\textbf{The named open problem: quantum guarded domain theory.} The
ingredients each exist --- guarded domain theory [Bir12], the completely
positive construction [Sel07] --- while their composition does not. Target:
presheaves on the ordered naturals valued in $\mathbf{CPM}$, $\guard$ the
shift, the guarded fixpoint the productive solution. Three walls: CPM is
not cartesian closed (the higher-order obstruction the quantum
$\lambda$-calculus's denotational program left open [SV06]); no infinite
stream lives in $\mathbf{CPM}(\mathbf{FdHilb})$, and the dagger structure
the physics needs is delicate under presheaf constructions; and guarded
fixpoint solvability wants $\omega$-limit enrichment the quantum side does
not yet have. Success would buy a denotation for the sigil, serialization
as a semantic fact, and the tick's scale-freeness as a theorem. The
construction is open.

\section{The tick and the boundary}

\dcn{knower_sustains_stream} \meas{}: one duplicable source sustains a
term of guarded-stream type --- one source, sustaining an infinite stream
notation of [GS], the seed
$\gknower : \tS \glto (\guard A \gfeed \gknower)$, instantiated $A := BH$
(horizon-delayed recursor; Fig.~\ref{fig:sigil}): the exterior of the loop
is a \emph{serialized} stream --- parallel interior process, step-indexed
exterior emission. The guard orders emission; it fixes no rate, temperature,
or spectrum, and no radiation is computed here. \interp{} The count is a theorem at the chain grain: emissions $=$
$\guard$-depth $=n$ at every stage (\dcn{chain_tick_buys_one} \meas{}),
and $n$ guarded ticks emit exactly $n$ while the spin emits exactly zero
(\dcn{guarded_vs_spin} \meas{}); emission-$\Rightarrow$-tick is not a
theorem (the extractor fires before any unfold) --- the count binds the
unfolding chain. The unfolding itself is anchored: each unfold consumes one
$\guard$ and is repeatable at every interval
(\dcn{fix_unfold_closed}, \dcn{fix_unfold_at_knower} \meas{}), and every
guarded agent step appends exactly one element
(\dcn{extractor_step_integrates} \meas{}). And the guard \emph{recedes}: the tower frames are priced
(zero budget) and stream is refuted at every one [CCS], with the base case
in-tree (\dcn{no_closed_stream} \meas{}: from nothing, nothing flows) ---
but a fix-rooted closed $\guard\,\mathrm{stream}$ is \emph{equivalent} to
the same question one guard higher, so no natural induction terminates, and
no monotone fix-sound price model of the known clause shape can close the
family: an impossibility theorem about the instruments [CCS]; the interior
is infinitely stratified for this question. \interp{} The fixpoint gate
(\texttt{AllNonLinear}) separates a duplicable ambient context from the
conserved linear stream it sustains; the counting model carries the stream's
budget on its ambient frame, and the arithmetic is a lemma at every scale $n$
(\dcn{holographic_budget} \meas{}): $n$ fires mint $n$ contents, the tower is
priced ${\ge}n$, and the knower source funding all of it is priced
$\emptyset$ yet occupies exactly one position, unchanged. Structurally --- and only
structurally: no metric, no field theory, no duality claim --- this is the
boundary/bulk separation of holography. \interp{}

\begin{figure}[t]
\centering
\begin{tikzpicture}
  \node[box] (src) at (0,0) {$S$};
  \node[box] (hrz) at (3,0) {$\guard BH$};
  \node[box] (out) at (6,0) {$\gknower$};
  \draw[wire] (src) -- node[lbl, above] {stream} (hrz);
  \draw[wire] (hrz) -- node[lbl, above] {one per tick} (out);
  \draw[wire] (hrz) to[out=120, in=60, looseness=1.6]
    node[lbl, above] {$\gfeed$ one interval per self-call} (hrz);
  \node[lbl, below=6pt of hrz] {the horizon-delayed recursor
    $\gknower : \tS \glto (\guard BH \gfeed \gknower)$};
\end{tikzpicture}
\caption{The seed of [GS] with the loop type instantiated as the reading
demands: a stream source feeding a horizon-delayed ($\guard$) recursor whose
feedback wire ($\gfeed$) pays one interval per self-call. The exterior sees
one emission per tick. \interp{It is a reading of a type, not a calculation
of radiation.}}
\label{fig:sigil}
\end{figure}

\section{The refutation: why the guard is load-bearing}

The machine-checked counterexample has a named anatomy \meas{}:
\dcn{the_cycle} (a duplicable source sustains a stream),
\dcn{the_swap} (the lete/pair exchange spends the content token as the
stream's next emission), \dcn{the_conduit} (a reader is synthesized to
receive it), \dcn{the_discard} (the reader is weakened away --- deleted).
Together: \dcn{unrestricted_counterexample} \meas{} --- a well-formed
derivation whose conclusion carries neither content nor record. Under the
dictionary, a \emph{non-unitary channel}: the state is not transformed, it is
spent into unreadability. \interp{} The kernel refuses it three times over
(\dcn{counterexample_not_noweak}: the discard ban alone suffices;
\dcn{counterexample_conduit_no_synthesis}: the synthesis ban suffices
independently; \dcn{counterexample_not_armed}: the full setup, whose armed
forms carry both laws, \dcn{no_contraction_armed} \meas{}) ---
redundant protection of a conserved quantity. \interp{} And the
two-worlds picture is one theorem: same generator, policy, start ---
guarded it halts and cannot diverge; unguarded it diverges
(\dcn{guard_halts_unguarded_spins}, \dcn{guarded_no_divergence}
\meas{}). And the unguarded
\emph{dynamics} never settle: \dcn{adversary_spin_diverges} and
\dcn{spin_pays_forever_buys_nothing} \meas{} --- cost forever, nothing
delivered, deterministically (from the spun state the only unguarded step is
the spin itself, \dcn{looped_only_spins} \meas{}; under the guarded
semantics a non-infinite run is finite --- silence persists,
\dcn{run_not_infinite_is_finite} \meas{}). One labeled sentence: \interp{a sector that permits the discard
admits, by theorem, a process that destroys content while its dynamics never
settle --- the \emph{shape} of a metastable phase; nothing here quantifies
it.} \textbf{The selection reading}: outside the kernel, loss is derived and
divergence is derived; inside, conservation is absolute --- a stable sector
and a lawful sector are the same sector, and no observer is needed --- its
ancestry is Wheeler's geometry-only foam and ``law without law'' [Whe55,
Whe57, Whe83], and, closest in spirit, cosmological natural selection,
which selects over parameters, never laws [Smo92, Smo06]; the anthropic
lineage [Car74] is explicitly not needed. Its one
tree-checkable fingerprint: the pre-lawful regime is \emph{lossy, not a
breeder} --- cloning stays forbidden even unrestricted (\dcn{no_contraction_content}
\meas{}); what fails outside is conservation, never non-duplication. The
shape has ancestry: any \emph{linear} non-unitary evolution (Hawking's
superscattering [Haw82, Page93]) and the no-broadcasting theorem's scope
[BCJFS96, BBLW07, Lin99] share it, and in type theory it is the affine
discipline [Wal04, Gir87]; ours is its packaging as theorems of one
rule-set. Structural ancestry, not identification: the foam lineage's loss
proposals (Hawking's superscattering [Haw76, Haw82], intrinsic decoherence
[EHNS84], gravitational reduction [Pen96]) fail unitarity globally or
wherever a black hole forms --- never outside a syntactically defined sector
--- and the modern verdict restores it [Page93, Haw05, Carlip23]; closest is
Cotler--Jensen's de Sitter evolution, unitary on a subspace [CJ19], a
dynamical, not grammatical, boundary --- and the cost argument of [BPS84]
shares its shape. \interp{}

\section{Experimental signatures}

A \proj{} throughout; physics inputs marked. (P1) The Bekenstein--Hawking
area law [Bek73, Haw75]: a radius-$R$ region is described by a screen of
$N \lesssim \pi R^2/l_p^2$ irreducible degrees of freedom. (P2) Assume the
tick spatialized --- bulk discretized at $L$ (an assumption the calculus does
not make). Then $(R/L)^3 \lesssim \pi R^2/l_p^2$ gives
\begin{equation}
L \;\gtrsim\; \pi^{-1/3}\,\bigl(l_p^2 R\bigr)^{1/3},
\qquad\text{i.e.}\qquad
L \approx 5\times10^{-15}\,\mathrm{m}\ \ (R = 4.4\times10^{26}\,\mathrm{m}):
\end{equation}
the observable bulk's voxel is nuclear-sized --- the MIP-map scaling, one bit
per screen pixel, and not ours: K\'arolyh\'azy's clock argument gave
$(\Delta s)^3 \gtrsim l_p^2 s$ in 1966 [Kar66]; Ng and van Dam re-derived it by the
cube-counting route [NgdV94, NgdV00, tHooft93, Suss95], and Ng's review runs our exact
derivation and prints the same order of magnitude for the $\alpha = 2/3$
branch [Ng03]. The
transverse variant $\sqrt{c\,t_P L}$ [Hog12] was tested by the Holometer's
paired 39-meter interferometers and is \emph{excluded} at $5.1\sigma$
(normalization below $44\%$ of the prediction at $95\%$ confidence)
[Chou16]; the volumetric family's channel is \emph{angular} --- wave-packet
widening blurring gamma-ray-burst images [Stein15], \emph{not} a timing
signature ($5\times10^{-15}\,$m is $1.7\times10^{-23}\,$s of light travel;
the picosecond figures are the transverse family's). The landscape is
contested, not clean: the 2003 coherence claims [Rag03, Lieu03] were
themselves corrected by the review we cite [Ng03] (the claimed exclusions
overestimated by $(L/\lambda)^{\alpha}$); the blazar program bounds
$\alpha \gtrsim 0.58$ (X-ray), $\alpha \gtrsim 0.67$ at GeV ---
``disfavoring the holographic model, but
perhaps not decisively'' [Perlman14] --- and $\gtrsim 0.72$ at TeV, just
above $2/3$; the GRB aggregation fits $\alpha = 0.650$--$0.735$ with
$2/3$ inside [Stein23]; and the field's methodology review names the two
$O(1)$-class freedoms (amplitude fixed only as $a \sim 1$ [Ng03]; distance
averaging unsettled) that keep every number provisional [NgP22]. The
projection stands where it stands: consistent with the GRB window,
strained by the TeV bound, undecidable until the averaging problem is
settled. And GRB 221009A's angular data are consistent with the $\alpha = 2/3$
model [Stein23], while the generalized K\'arolyh\'azy model has just been
excluded on the neighboring gravity-induced-decoherence channel ---
$R_K > 4.64\,$m at $95\%$ confidence against a predicted $R_K < 1.98\,$m
[Bort26] --- not the blurring channel, and not the gravitational-wave
band, which points the wrong instruments at this scale.
The $1/3$ exponent is fixed by the entropy bound's dimensionality, not by
the guard. \textbf{The guard fixes no scale.} We looked for a second
handle and state its absence: the time-domain version of the same scaling
gives $\delta t \approx 1.7\times10^{-23}\,$s over a cosmological baseline
--- ten orders below any timing instrument; the angular channel is the
only live one. A detection would support the
projection; a null result constrains the projection family; neither touches
the theorems.

\section{Conclusion}

The two conservation facts appear in this grammar as stated: cloning
forbidden absolutely, conservation absolute in the discard-free sector, loss
derived and divergence derived outside it, the exterior serialized by the
tick, the boundary structural. The one projection the grammar motivates comes
with its falsifiers named and its epicycles refused. In this reading, physics
is, possibly, the study of the well-typed sector --- the paper's largest
\interp{} and its smallest claim.

\section*{Scope and limits}
The bridge is a mirror, not a derivation: the semantic-bridge construction above builds a functor from
the finitary fragment into $\mathbf{CPM}(\mathbf{FdHilb})$ that preserves
processes and mirrors the two laws categorically, but functors do not
reflect non-derivability --- the grammar's no-go theorems are stronger than
their denotations --- and no bridge to spacetimes or to the guarded
fragment is built. The correspondence otherwise remains a dictionary, not
an isomorphism. The counting model is natural-number-valued (no inner
product, no amplitudes). The calculus contains no geometry; all scales in
the experimental section are physics inputs to a projection. Every
\meas{voucher} is a named \dcn{} declaration verified in-tree by this
artifact's \texttt{scripts/sync-check.sh}.

\begin{thebibliography}{99}\small
\bibitem{WZ82} Wootters, W.\,K., \& Zurek, W.\,H. (1982). A single quantum
  cannot be cloned. \emph{Nature} 299, 802--803.
\bibitem{PB00} Pati, A.\,K., \& Braunstein, S.\,L. (2000). Impossibility of
  deleting an unknown quantum state. \emph{Nature} 404, 164--165.
\bibitem{AC04} Abramsky, S., \& Coecke, B. (2004). A categorical semantics of
  quantum protocols. \emph{Proc.\ LICS 2004}, 415--425.
\bibitem{Sel07} Selinger, P. (2007). Dagger compact closed categories and
  completely positive maps. \emph{Electronic Notes in Theoretical Computer
  Science} 170, 139--163.
\bibitem{Wig59} Wigner, E.\,P. (1959). \emph{Group Theory and Its Application
  to the Quantum Mechanics of Atomic Spectra.} Academic Press.
\bibitem{Gir87} Girard, J.-Y. (1987). Linear logic. \emph{Theoretical
  Computer Science} 50(1), 1--101.
\bibitem{Wal04} Walker, D. (2004). Substructural type systems. In B.\,C.\
  Pierce (ed.), \emph{Advanced Topics in Types and Programming Languages},
  MIT Press, 3--44.
\bibitem{Abr09} Abramsky, S. (2009). No-cloning in categorical quantum
  mechanics. In \emph{Semantic Techniques in Quantum Computation}, CUP, 1--28.
\bibitem{SV06} Selinger, P., \& Valiron, B. (2006). A lambda calculus for
  quantum computation with classical control. \emph{Math.\ Struct.\ in
  Comp.\ Sci.} 16(3), 527--552.
\bibitem{PRZ17} Paykin, J., Rand, R., \& Zdancewic, S. (2017). QWIRE: a core
  language for quantum circuits. \emph{Proc.\ POPL 2017}, PACMPL 1(POPL).
\bibitem{BHL20} Bordg, A., Lachnitt, H., \& He, Y. (2020). Certified quantum
  computation in Isabelle/HOL. \emph{J.\ Automated Reasoning} 65, 691--709.
\bibitem{ILV26} Lamiroy, J., Valiron, B., \& Vilmart, R. (2026).
  Intuitionistic unitary linear logic. arXiv:2609.02661.
\bibitem{Lin99} Lindblad, G. (1999). A general no-cloning theorem. \emph{
  Letters in Math.\ Phys.} 47, 189--196.
\bibitem{BCJFS96} Barnum, H., Caves, C.\,M., Fuchs, C.\,A., Jozsa, R., \&
  Schumacher, B. (1996). Noncommuting mixed states cannot be broadcast.
  \emph{Phys.\ Rev.\ Lett.} 76, 2818--2821.
\bibitem{BBLW07} Barnum, H., Barrett, J., Leifer, M., \& Wilce, A. (2007).
  A generalized no-broadcasting theorem. \emph{Phys.\ Rev.\ Lett.} 99, 240501.
\bibitem{HHSS05} Horodecki, M., Horodecki, R., Sen(De), A., \& Sen, U.
  (2005). Common origin of no-cloning and no-deleting principles:
  conservation of information. \emph{Found.\ Phys.} 35, 2041--2049.
\bibitem{Nak00} Nakano, H. (2000). A modality for recursion. \emph{Proc.\
  LICS 2000}.
\bibitem{AM13} Atkey, R., \& McBride, C. (2013). Productive coprogramming
  with guarded recursion. \emph{Proc.\ ICFP 2013}.
\bibitem{Bir12} Birkedal, L., M\o gelberg, K.\,E., Schwinghammer, J., \&
  St\o vring, K. (2012). First steps in synthetic guarded domain theory.
  \emph{LMCS} 8(4).
\bibitem{Whe55} Wheeler, J.\,A. (1955). Geons. \emph{Phys.\ Rev.} 97, 511--536.
\bibitem{Whe57} Wheeler, J.\,A. (1957). On the nature of quantum
  geometrodynamics. \emph{Ann.\ Phys.} 2, 604--614.
\bibitem{Whe83} Wheeler, J.\,A. (1983). Law without law. In Wheeler \& Zurek
  (eds.), \emph{Quantum Theory and Measurement}, Princeton UP, 182--213.
\bibitem{Haw76} Hawking, S.\,W. (1976). Breakdown of predictability in
  gravitational collapse. \emph{Phys.\ Rev.\ D} 14, 2460--2473.
\bibitem{Haw82} Hawking, S.\,W. (1982). The unpredictability of quantum
  gravity. \emph{Commun.\ Math.\ Phys.} 87, 395--415.
\bibitem{EHNS84} Ellis, J., Hagelin, J.\,S., Nanopoulos, D.\,V., \&
  Srednicki, M. (1984). Search for violations of quantum mechanics.
  \emph{Nucl.\ Phys.\ B} 241, 381--405.
\bibitem{Pen96} Penrose, R. (1996). On gravity's role in quantum state
  reduction. \emph{Gen.\ Rel.\ Grav.} 28, 581--600.
\bibitem{Page93} Page, D.\,N. (1993). Information in black hole radiation.
  \emph{Phys.\ Rev.\ Lett.} 71, 3743--3746.
\bibitem{Haw05} Hawking, S.\,W. (2005). Information loss in black holes.
  \emph{Phys.\ Rev.\ D} 72, 084013.
\bibitem{Carlip23} Carlip, S. (2023). Spacetime foam: a review. \emph{Rep.\
  Prog.\ Phys.} 86, 066001.
\bibitem{BPS84} Banks, T., Susskind, L., \& Peskin, M.\,E. (1984).
  Difficulties for the evolution of pure states into mixed states.
  \emph{Nucl.\ Phys.\ B} 244, 125--134.
\bibitem{CJ19} Cotler, J., \& Jensen, K. (2021). Emergent unitarity in de
  Sitter from matrix integrals. \emph{JHEP} 12, 089.
\bibitem{Smo92} Smolin, L. (1992). Did the universe evolve? \emph{Class.\
  Quantum Grav.} 9, 173--192.
\bibitem{Smo06} Smolin, L. (2006). The status of cosmological natural
  selection. arXiv:hep-th/0612185.
\bibitem{Car74} Carter, B. (1974). Large number coincidences and the
  anthropic principle in cosmology. \emph{IAU Symposium 63}, 291--298.
\bibitem{Bek73} Bekenstein, J.\,D. (1973). Black holes and entropy.
  \emph{Phys.\ Rev.\ D} 7, 2333--2346.
\bibitem{Haw75} Hawking, S.\,W. (1975). Particle creation by black holes.
  \emph{Commun.\ Math.\ Phys.} 43, 199--220.
\bibitem{tHooft93} 't Hooft, G. (1993). Dimensional reduction in quantum
  gravity. \emph{Conf.\ Proc.\ C} 930308, 284; arXiv:gr-qc/9310026.
\bibitem{Suss95} Susskind, L. (1995). The world as a hologram. \emph{J.\ Math.\
  Phys.} 36, 6377--6396.
\bibitem{Kar66} K\'arolyh\'azy, F. (1966). Gravitation and quantum mechanics
  of macroscopic objects. \emph{Nuovo Cimento A} 42, 390.
\bibitem{NgdV94} Ng, Y.\,J., \& van Dam, H. (1994). Limit to space-time
  measurement. \emph{Mod.\ Phys.\ Lett.\ A} 9, 335.
\bibitem{NgdV00} Ng, Y.\,J., \& van Dam, H. (2000). Measuring the foaminess
  of space-time with gravity-wave interferometers. \emph{Found.\ Phys.} 30, 795.
\bibitem{Ng03} Ng, Y.\,J. (2003). Selected topics in Planck-scale physics.
  \emph{Mod.\ Phys.\ Lett.\ A} 18, 1073.
\bibitem{Hog12} Hogan, C.\,J. (2012). Interferometers as probes of Planckian
  quantum geometry. \emph{Phys.\ Rev.\ D} 85, 064007.
\bibitem{Chou16} Chou, A.\,S., et al. (2016). First measurements of high
  frequency cross-spectra from a pair of large Michelson interferometers.
  \emph{Phys.\ Rev.\ Lett.} 117, 111102.
\bibitem{Stein15} Steinbring, E. (2015). Detectability of Planck-scale-induced
  blurring with gamma-ray bursts. \emph{Astrophys.\ J.} 802, 38.
\bibitem{Rag03} Ragazzoni, R., Turatto, M., \& Gaessler, W. (2003). The lack
  of observational evidence for the quantum structure of spacetime at Planck
  scales. \emph{Astrophys.\ J.\ Lett.} 587, L1--L4.
\bibitem{Perlman14} Perlman, E., Ng, Y.\,J., et al. (2014). New constraints
  on quantum foam models from X-ray and gamma-ray observations.
  arXiv:1411.7262; proc.\ 14th Marcel Grossmann Meeting (2017), 3935--3941.
\bibitem{NgP22} Ng, Y.\,J., \& Perlman, E. (2022). Probing spacetime foam
  with extragalactic sources of high-energy photons. \emph{Universe} 8(7),
  382.
\bibitem{Lieu03} Lieu, R., \& Hillman, L.\,W. (2003). The phase coherence of
  light from extragalactic sources: direct evidence against first-order
  Planck-scale fluctuations in time and space. \emph{Astrophys.\ J.\ Lett.}
  585, L77--L80.
\bibitem{Stein23} Steinbring, E. (2023). Holographic quantum-foam blurring is
  consistent with observations of gamma-ray burst GRB 221009A.
  \emph{Galaxies} 11, 115.
\bibitem{Bort26} Bortolotti, N., et al. (2026). Experimental exclusion of a
  generalized K\'arolyh\'azy gravity-induced decoherence model. \emph{New
  J.\ Phys.} 28, 064511.
\bibitem{CCS} \emph{The Care Calculus: A Mechanized Linear Type Theory for
  Safe Agents.} Companion mechanization and paper (in-tree dcare artifact).
\bibitem{GS} \emph{The Guarded Syntax: A Notation for Linear Resource
  Calculi.} Program artifact (in-tree notation artifact).
\bibitem{GL} \emph{The Guarded Line: What Metacircular Proof Theory and
  Self-Model Dynamics Share.} Program artifact (in-tree isomorphism
  artifact).
\end{thebibliography}

\end{document}
